\documentclass{article}
\usepackage{PRIMEarxiv} % Core package for the PrimeAI Template
\usepackage{amsmath, amssymb, amsthm, bm, dcolumn} % Add amsthm here for the proof environment
\usepackage[numbers,sort&compress]{natbib} % Natbib for citations
\usepackage{graphicx} % For high-quality images
\usepackage{multicol} % Multi-column figures/tables
\usepackage{hyperref} % Hyperlinks
\usepackage{listings} % Code listings
\usepackage{authblk} % For structured affiliations
% Define theorem style (optional)
\theoremstyle{plain}
\newtheorem{theorem}{Theorem}
\renewcommand{\qedsymbol}{}

% Custom commands for primes and qubit encoding
\newcommand{\primeQubit}[1]{|p_{#1}\rangle}
\newcommand{\primeGate}[1]{U_{p_{#1}}}

\usepackage{wrapfig}
\usepackage[pscoord]{eso-pic}
\usepackage[fulladjust]{marginnote}
\reversemarginpar

% Typesetting improvements without footnote patching
\usepackage[protrusion=true, expansion=true, tracking=false]{microtype}
\microtypecontext{spacing=nonfrench}

% Line numbers
\usepackage[right]{lineno}

% Text layout - adjust as needed
\raggedright
\setlength{\parindent}{0.5cm}
\textwidth 5.25in 
\textheight 8.75in

% Set double spacing
\usepackage{setspace} 
\doublespacing

% Adjust width for specific content
\usepackage{changepage}

% Adjust caption style
\usepackage[aboveskip=1pt,labelfont=bf,labelsep=period,singlelinecheck=off]{caption}

% Remove brackets from references
\makeatletter
\renewcommand{\@biblabel}[1]{\quad#1.}
\makeatother

% Header, footer, and page numbers
\usepackage{lastpage,fancyhdr}
\pagestyle{fancy}
\fancyhf{}
\fancyhead[L]{Preprint - PrimeAI Enhanced Template}
\fancyfoot[C]{\scriptsize Multiplicity Theory © 2024 Ryan Van Gelder - Citizen Gardens \\ Licensed Under MIT and CC BY-NC-SA 4.0.}
\fancyfoot[R]{Page \thepage\ of \pageref{LastPage}}
\renewcommand{\footrule}{\hrule height 2pt \vspace{2mm}}



\title{The Implications of Multiplicity Theory \\ in Machine Learning}

\author{Ryan O. Van Gelder \& Nicholas Galioto\\
Citizen Gardens - The Foundation of Multiplicity \\
\texttt{info@citizengardens.org}}

\begin{document}
\maketitle

\begin{abstract}
Multiplicity Theory introduces innovative computational frameworks rooted in prime-based encoding, recursive feedback, tensor networks, and quantum-inspired dynamics. This article provides a comprehensive overview of the theory's implications for machine learning and its mathematical foundations. Applications range from natural language processing and computer vision to optimization and reinforcement learning, demonstrating improved efficiency, scalability, and adaptability.
\end{abstract}

\section{Introduction}
Multiplicity Theory redefines computational paradigms by integrating prime number theory, tensor networks, and recursive dynamics. In machine learning, it enhances robustness, adaptability, and scalability. This paper outlines its implications and the mathematical frameworks underlying these advancements.




\nocite{*}
\bibliographystyle{unsrt}
\bibliography{references}

\end{document}